% from main.tex
%%%%%%%%%%  scan idx 404  |  folio 391  %%%%%%%%%%
pour tout \(i\) un \(\mathcal O_{X_i}\)-Module localement libre de type fini \(F_i\) en posant~:
\[
  X_k=\mathbb P(F_{k-1});\qquad
  F_k=\Ker\bigl(f_k^*(F_{k-1})\longrightarrow \mathcal O_{X_k}(1)\bigr),
\]
où \(f_k\) est le morphisme structural \(X_k\to X_{k-1}\). On voit alors qu'il suffit de montrer la proposition dans le cas où \(X\) est un fibré projectif sur \(S\), le cas d'un fibré de drapeaux s'en déduisant par récurrence. Or ce cas résulte immédiatement de 1.1 (resp. 2.1).

\textbf{Théorème 3.2.} Soit \(S\) un schéma quasi-compact. Le pré-\(\lambda\)-anneau \(K^{\bullet}(S)\), défini à partir des \(\mathcal O_S\)-Modules localement libres de type fini, est un \(\lambda\)-anneau~(V 2.4).

L'homomorphisme \(\lambda_t:K^{\bullet}(S)\longrightarrow 1+K^{\bullet}(S)[[t]]^+\) étant additif, il suffit, pour vérifier que c'est un \(\lambda\)-homomorphisme, de faire la vérification pour des éléments du type \(\cl(E)\), où \(E\) est un \(\mathcal O_S\)-Module localement libre de type fini. On est donc ramené à montrer que pour tout couple d'éléments \(E,F\) de \(K^{\bullet}(S)\), classes de \(\mathcal O_S\)-Modules localement libres de type fini \(E,F\), les relations V 2.4.2 sont vérifiées :
\[
  \lambda^i(EF)=P_i\bigl(\lambda^1(E),\ldots,\lambda^i(E);\lambda^1(F),\ldots,\lambda^i(F)\bigr);
\]
\[
  \lambda^j(\lambda^i(E))=Q_{i,j}\bigl(\lambda^1(E),\ldots,\lambda^{ij}(E)\bigr).
\]
Comme \(S\) est un quasi-compact, on peut trouver un recouvrement fini de \(S\) par des ouverts disjoints tel que sur chacun \(E\) et \(F\) soient de rang constant. Or le \(K^{\bullet}\) d'une somme directe finie d'ouverts est le produit des \(K^{\bullet}\) de ces ouverts. Il suffit donc de vérifier les relations V 2.4.2 dans chacun des \(K^{\bullet}\) des ouverts du recouvrement. On peut donc supposer \(E\) et \(F\) de rang constant sur \(S\).

